\section{Results}
\label{sec:results}

This section presents EPGSpeller results as a single narrative organized around the protocol map in Table~\ref{tab:protocol_map}.
We first compare within-user conditions with user calibration against cross-user transfer using the main vector/GRU recognizer and greedy string decoding before lexicon projection.
We then connect to SilentSpeller-style unseen-word evaluation, evaluate low-shot calibration, multi-source transfer, and preprocessing controls, and summarize engineering analyses.

\subsection{Main Results: Strong Within-User Performance, Large Cross-User Degradation}
\label{sec:results_main}

Table~\ref{tab:main_protocols} summarizes the main results.
With the same vector/GRU recognizer and greedy string decoding, there is a large gap between within-user conditions with target-user calibration and cross-user transfer without it.
P1 (word holdout, within user) reaches \cerstats{0.1796}{0.0844} CER (mean \textpm\ std over four random seeds).
P2 (instance holdout, within user) reaches \cerstats{0.1451}{0.0627} CER (three seeds).
In contrast, P3 (cross-user transfer) degrades to \cerstats{0.6909}{0.1332} CER (six transfer directions).
The P1/P2 mean \textpm\ std values are aggregated over random-seed repeats (n=4 and n=3, respectively); P3 is aggregated over transfer directions (n=6).
Lexicon-projected values (LEX) are reported as a supplementary diagnostic for closed-vocabulary conditions.
LEX further reduces error under within-user conditions, but the lexicon-free claim rests on CER before lexicon projection.

\begin{table}[!b]
\centering
\caption{Main results for the vector/GRU recognizer (mean \textpm\ std). n = seeds (P1/P2) or directions (P3). LEX: lexicon-projected CER; RTF: real-time factor.}
\label{tab:main_protocols}
\setlength{\tabcolsep}{2pt}
\begin{tabular}{@{}llll@{}}
\toprule
Cond. & CER & LEX & RTF \\
\midrule
P1 & \cerstats{0.1796}{0.0844} & \cerstats{0.1018}{0.0679} & \cerstats{0.000628}{0.000232} \\
P2 & \cerstats{0.1451}{0.0627} & \cerstats{0.0747}{0.0476} & \cerstats{0.000107}{0.000010} \\
P3 & \cerstats{0.6909}{0.1332} & \cerstats{0.6442}{0.0604} & \cerstats{0.000110}{0.000014} \\
\bottomrule
\end{tabular}
\end{table}

\subsection{SilentSpeller Bridge and Low-Shot Calibration}

\begin{table}[t]
\centering
\caption{Low-shot calibration and evaluability on p3 (CER mean \textpm\ std).}
\label{tab:kshot_curve}
\setlength{\tabcolsep}{3pt}
\begin{tabular}{@{}lll@{}}
\toprule
Condition & CER & n \\
\midrule
Zero-shot, single source & \cerstats{0.8403}{0.0948} & 8 \\
Zero-shot, multi-source & \cerstats{0.8381}{0.0336} & 4 \\
k=1 & \cerstats{0.3503}{0.0600} & 4 \\
k=2 & \cerstats{0.2458}{0.0085} & 4 \\
k=4 & --- (14 words) & --- \\
k=8 & infeasible & 0 \\
\bottomrule
\end{tabular}
\end{table}

We assess an auxiliary SilentSpeller bridge on participant p1 with 100 held-out test words from the same release.
The vector/GRU recognizer reaches \cerstats{0.1133}{0.0128} CER averaged over four seeds under this non-interactive setup.
This sanity check confirms unseen-word decoding within the release inventory; it does not compare against SilentSpeller-era interactive decoders (deferred; Section~\ref{sec:evaluation}).

Low-shot calibration experiments quantify cross-user transfer as a calibration-quantity problem.
On p3, zero-shot single-source transfer remains at \cerstats{0.8403}{0.0948} CER (n=8), and zero-shot multi-source transfer at \cerstats{0.8381}{0.0336} CER (n=4).
One example per word improves to \cerstats{0.3503}{0.0600} CER (n=4); two examples per word to \cerstats{0.2458}{0.0085} CER (n=4).
Higher-$k$ evaluation is constrained by public data density (Table~\ref{tab:calibration_setup}).

\subsection{Transfer Controls and Engineering Analyses}

\begin{table}[t]
\centering
\caption{Engineering analyses: spatial front ends and electrode reduction (greedy CER, mean \textpm\ std).}
\label{tab:spatial_electrode}
\setlength{\tabcolsep}{2pt}
\begin{minipage}[t]{0.47\columnwidth}
\centering
\textbf{(a) P1 spatial front end}\par
\begin{tabular}{@{}ll@{}}
\toprule
Front end & CER \\
\midrule
vector/GRU & \cerstats{0.18}{0.08} \\
row-column & \cerstats{0.20}{0.09} \\
grid & \cerstats{0.31}{0.15} \\
patch & \cerstats{0.30}{0.07} \\
\bottomrule
\end{tabular}
\end{minipage}
\hfill
\begin{minipage}[t]{0.47\columnwidth}
\centering
\textbf{(b) K=64 electrode reduction}\par
\begin{tabular}{@{}lll@{}}
\toprule
Protocol & Method & CER \\
\midrule
P1 & top-k & \cerstats{0.195}{0.083} \\
P1 & spatial & \cerstats{0.191}{0.078} \\
P2 & top-k & \cerstats{0.154}{0.059} \\
P2 & spatial & \cerstats{0.154}{0.062} \\
\bottomrule
\end{tabular}
\end{minipage}
\end{table}

Transfer controls test whether multi-source training or simple preprocessing can close the cross-user gap without target-user calibration.
Multi-source transfer partially improves cross-user performance but does not resolve the calibration problem.
When p4 is the target user, single-source cross-user transfer with all training sources to p4 (P3) reaches 0.686 CER, whereas multi-source transfer (P3MS) reaches 0.582 CER---both remain far above the within-user P1/P2 range in Table~\ref{tab:main_protocols} (roughly 0.15--0.18 CER).
Preprocessing controls test whether cross-user failure can be explained by simple scale differences or contact-rate mismatch alone.
A sweep over raw input, standardization, contact-rate matching, target-neighborhood distribution normalization, and PCA-based standardization leaves best protocol averages at 0.6426 for P3 and 0.5692 for P3MS.

We analyze how SmartPalate signals should be represented and how future devices might be simplified.
On P1, the vector/GRU path reaches \cerstats{0.18}{0.08} CER; row-column pooling is close at \cerstats{0.20}{0.09}, whereas grid and patch variants degrade to roughly \cerstats{0.31}{0.15} and \cerstats{0.30}{0.07} (four seeds; Table~\ref{tab:spatial_electrode}(a)).
Two-dimensional grid assumptions are therefore not consistently advantageous on sparse, user-specific contacts.

Electrode reduction at K=64 channels retains near-main P1/P2 performance across top-k, spatial, transfer-ranked, and random selection (Table~\ref{tab:spatial_electrode}(b)).
Augmentation suppresses degradation from local missing contacts and channel dropout.
Taken together, the spatial-front-end, electrode-reduction, and preprocessing sweeps bound representation and normalization choices but leave the P3 calibration gap essentially unchanged relative to within-user P1/P2 performance.
None of these engineering analyses closes the P3 cross-user gap without target-user examples.
